\section{Conclusion}
\label{sec:conclusion}

The evidence reviewed here raises questions about the
assumption that efficiency gains automatically lead to
CO\textsubscript{2} reductions in AI deployment. Both the
direct and indirect cases examined follow a consistent logic:
lower costs expand demand, and resource consumption grows
faster than efficiency improvements alone can offset.

From a research perspective, the most pressing gap concerns the
indirect channel. While direct rebound effects -- rising GPU
shipments and data center energy demand -- are increasingly
documented, indirect mechanisms remain poorly quantified at the
macro level. Future work should examine sector-specific induced
demand, especially in logistics and e-commerce, using
longitudinal cross-national data and standardized emissions
accounting. Methodological advances in attributing consumption
growth to AI-specific cost reductions would substantially
strengthen the evidence base.

For practice, voluntary efficiency improvements may be
insufficient in isolation. Realizing the ecological potential
of AI may require policy intervention: binding resource caps,
mandatory environmental reporting, and growth targets that do
not presuppose the indefinite expansion of AI infrastructure.
